\subsection{Persistent components and glass constraints}

\begin{theorem}[Persistent \(H^0\) components only merge]\label{thm:persistent-h0-components-merge}\lean{Obsmetrologia.Calculus.Order.components_only_merge}\leanok\uses{thm:cover-nerve}
In a graph filtration, path-connected components are monotone: increasing the
graph can merge components but cannot split an existing connection.
\end{theorem}

\begin{theorem}[Kinetically constrained activation is closure]\label{thm:kcm-activation-closure}\lean{Obsmetrologia.AperiodicOrder.GlassDynamics.activation_is_facilitation_closure}\leanok\uses{thm:causal-settling-transitive-closure}
A kinetically constrained activation process equals the inductive facilitation
closure generated by the seed and the facilitation graph.
\end{theorem}

\begin{theorem}[Finite glass overlap is symmetric]\label{thm:glass-overlap-symmetric}\lean{Obsmetrologia.AperiodicOrder.GlassDynamics.overlapCount_symm}\leanok\uses{thm:valuation}
The finite replica-overlap count counts the sites where two configurations
agree.  The count is symmetric, and self-overlap fills the whole finite site
list in the companion theorem.
\end{theorem}

\begin{theorem}[Zero mobility is agreement on support]\label{thm:glass-mobility-zero-agreement}\lean{Obsmetrologia.AperiodicOrder.GlassDynamics.mobility_zero_iff_agree_on_support}\leanok\uses{thm:glass-overlap-symmetric}
The dynamic-heterogeneity count is the finite mobility count.  It vanishes
exactly when two configurations agree on every listed site.
\end{theorem}

\begin{theorem}[Basin descent has no cycle]\label{thm:basin-no-descending-cycle}\lean{Obsmetrologia.AperiodicOrder.GlassDynamics.BasinGraph.no_descending_cycle}\leanok\uses{thm:no-periodic-orbits}
A finite basin graph with strictly decreasing energy along descent has no
descending cycle. This is the finite landscape H-theorem.
\end{theorem}

\begin{proposition}[Glass substance bundle]\label{thm:glass-substance-bundle}\lean{Obsmetrologia.AperiodicOrder.Crystal.Glass}\leanok\uses{thm:charge-coloured-sofic, thm:lenticule-hole, thm:lenticle-identification, thm:percolation-threshold-law, thm:homometric-finite-diffraction-mass, thm:kcm-activation-closure, thm:basin-no-descending-cycle, thm:finite-thermodynamic-transition-production}
A glass substance is the bundled material carrier whose layers are
Tiles, Beads, Crystals, Quasicrystals, and the finite glass diagnostics:
condensed bead placement, configuration overlap, mobility, basin descent,
kinetic activation, and the stationary temperature/Landauer transition.
\end{proposition}

\begin{theorem}[Beads sit between Tiles and Crystals]\label{thm:glass-beads-between-tiles-and-crystals}\lean{Obsmetrologia.AperiodicOrder.Crystal.Glass.beads_between_tiles_and_crystals}\leanok\uses{thm:glass-substance-bundle}
The glass substance projection from crystals back down to tiles passes through
the bead layer.  This records Beads as a genuine intermediate carrier, not a
prose synonym for either tiles or crystals.
\end{theorem}

\begin{theorem}[Beads realize lenticles as clopen overlaps]\label{thm:glass-bead-material-realization}\lean{Obsmetrologia.AperiodicOrder.Crystal.BeadMaterialRealization, Obsmetrologia.AperiodicOrder.Crystal.BeadMaterialRealization.lenticle_eq_overlap, Obsmetrologia.AperiodicOrder.Crystal.BeadMaterialRealization.lenticle_sub_left, Obsmetrologia.AperiodicOrder.Crystal.BeadMaterialRealization.lenticle_sub_right, Obsmetrologia.AperiodicOrder.Crystal.BeadMaterialRealization.covers_lenticle_pair, Obsmetrologia.AperiodicOrder.Crystal.Beads.geometricLenticle, Obsmetrologia.AperiodicOrder.Crystal.Beads.geometricLenticle_eq_overlap, Obsmetrologia.AperiodicOrder.Crystal.Beads.geometricLenticle_sub_left, Obsmetrologia.AperiodicOrder.Crystal.Beads.geometricLenticle_sub_right, Obsmetrologia.AperiodicOrder.Crystal.Beads.covers_geometricLenticle_pair, Obsmetrologia.AperiodicOrder.Crystal.Crystals.beadRealization, Obsmetrologia.AperiodicOrder.Crystal.Crystals.beadGeometricLenticle_eq_overlap, Obsmetrologia.AperiodicOrder.Crystal.Crystals.covers_beadGeometricLenticle_pair, Obsmetrologia.AperiodicOrder.Crystal.Quasicrystals.beadRealization, Obsmetrologia.AperiodicOrder.Crystal.Quasicrystals.beadGeometricLenticle_eq_overlap, Obsmetrologia.AperiodicOrder.Crystal.Quasicrystals.covers_beadGeometricLenticle_pair, Obsmetrologia.AperiodicOrder.Crystal.Glass.beadRealization_eq_crystal_beadRealization, Obsmetrologia.AperiodicOrder.Crystal.Glass.beadRealization_eq_quasicrystal_beadRealization, Obsmetrologia.AperiodicOrder.Crystal.Glass.beadGeometricLenticle_eq_overlap, Obsmetrologia.AperiodicOrder.Crystal.Glass.covers_beadGeometricLenticle_pair}\leanok\uses{thm:glass-beads-between-tiles-and-crystals, thm:lenticle-identification}
The bead layer carries a finite material realization: two clopen material
patches cover an ambient patch, and the bead lenticle is their overlap.  Pulling
the cover back to the lenticle gives the duplicated overlap cover.  Crystals
and quasicrystals retain this site-level realization unchanged.  This is a
finite clopen geometry result, not an analytic embedding or continuum mechanics
claim.
\end{theorem}

\begin{theorem}[Beads carry the port law into Crystals]\label{thm:glass-beads-port-law-condensate}\lean{Obsmetrologia.AperiodicOrder.Crystal.Beads.thread_pass_ports_distinct, Obsmetrologia.AperiodicOrder.Crystal.Beads.over_occupation_inconsistency, Obsmetrologia.AperiodicOrder.Crystal.Beads.mirror_preserves_condensate_inconsistency, Obsmetrologia.AperiodicOrder.Crystal.Crystals.beadThreadLegal, Obsmetrologia.AperiodicOrder.Crystal.Crystals.bead_pass_ports_distinct, Obsmetrologia.AperiodicOrder.Crystal.Crystals.bead_over_occupation_inconsistency, Obsmetrologia.AperiodicOrder.Crystal.Crystals.mirror_preserves_beadCondensate_inconsistency, Obsmetrologia.AperiodicOrder.Crystal.Glass.crystal_beadThreadLegal, Obsmetrologia.AperiodicOrder.Crystal.Glass.crystal_bead_pass_ports_distinct, Obsmetrologia.AperiodicOrder.Crystal.Glass.beadCondensate_eq_crystal_beadCondensate}\leanok\uses{thm:glass-beads-between-tiles-and-crystals, thm:lenticule-hole}
The bead layer is not only an intermediate name.  It carries the lenticular
thread law: every retained thread pass enters and exits through distinct ports.
It also carries the condensate charge of bead-thread occupation, so over-
occupation is the held contradiction pole and mirroring preserves that
contradiction content.  Crystals and glass substances project these bead laws
upward unchanged.
\end{theorem}

\begin{theorem}[Bead multi-adjacency carries into Crystals]\label{thm:glass-bead-multi-adjacency-carrier}\lean{Obsmetrologia.AperiodicOrder.Crystal.BeadAdjacency, Obsmetrologia.AperiodicOrder.Crystal.BeadAdjacency.tuple_eq, Obsmetrologia.AperiodicOrder.Crystal.Beads.adjacencyTuple_eq, Obsmetrologia.AperiodicOrder.Crystal.Crystals.beadAdjacency, Obsmetrologia.AperiodicOrder.Crystal.Crystals.beadAdjacency_tuple_eq, Obsmetrologia.AperiodicOrder.Crystal.Quasicrystals.beadAdjacency, Obsmetrologia.AperiodicOrder.Crystal.Quasicrystals.beadAdjacency_tuple_eq, Obsmetrologia.AperiodicOrder.Crystal.Glass.beadAdjacency_eq_crystal_beadAdjacency, Obsmetrologia.AperiodicOrder.Crystal.Glass.beadAdjacency_eq_quasicrystal_beadAdjacency}\leanok\uses{thm:glass-beads-port-law-condensate, thm:bead-four-adjacency-carrier}
The bead layer carries four named adjacency graphs on the same lenticule
carrier: thread, contact, cell, and visible.  Crystals and quasicrystals retain
that adjacency bundle unchanged.  This is the carrier proved for
multi-adjacency; pairwise separation of the four relations and the visible
sofic-factor theorem remain frontier claims.
\end{theorem}

\begin{theorem}[Quasicrystals retain the bead condensate]\label{thm:glass-quasicrystal-retains-bead-condensate}\lean{Obsmetrologia.AperiodicOrder.Crystal.Quasicrystals.beadPassCount, Obsmetrologia.AperiodicOrder.Crystal.Quasicrystals.beadCondensate, Obsmetrologia.AperiodicOrder.Crystal.Quasicrystals.beadThreadLegal, Obsmetrologia.AperiodicOrder.Crystal.Quasicrystals.bead_pass_ports_distinct, Obsmetrologia.AperiodicOrder.Crystal.Quasicrystals.bead_over_occupation_inconsistency, Obsmetrologia.AperiodicOrder.Crystal.Quasicrystals.mirror_preserves_beadCondensate_inconsistency, Obsmetrologia.AperiodicOrder.Crystal.Quasicrystals.finite_diffraction_mass_univ, Obsmetrologia.AperiodicOrder.Crystal.Glass.beadCondensate_eq_quasicrystal_beadCondensate, Obsmetrologia.AperiodicOrder.Crystal.Glass.quasicrystal_bead_over_occupation_inconsistency, Obsmetrologia.AperiodicOrder.Crystal.Glass.quasicrystal_mirror_preserves_beadCondensate_inconsistency}\leanok\uses{thm:glass-beads-port-law-condensate, thm:finite-diffraction-mass}
The quasicrystal layer does not erase the bead layer.  It carries the same
bead-thread legality, bead condensate charge, over-occupation inconsistency, and
mirror-invariant contradiction content while adding the finite diffraction
readout of the crystal occupation pattern.  This is the bridge proved from
lenticular bead support to quasicrystal finite diffraction mass; the analytic
diffraction measure remains a frontier claim elsewhere.
\end{theorem}

\begin{theorem}[Bead placement carries percolation into Quasicrystals]\label{thm:glass-bead-placement-percolation}\lean{Obsmetrologia.AperiodicOrder.Crystal.Beads.condensedContacts, Obsmetrologia.AperiodicOrder.Crystal.Beads.placement_cells_eq_condensed, Obsmetrologia.AperiodicOrder.Crystal.Beads.loadBearing_of_condensed_crosses_threshold, Obsmetrologia.AperiodicOrder.Crystal.Crystals.beadPlacement, Obsmetrologia.AperiodicOrder.Crystal.Crystals.beadCondensedContacts, Obsmetrologia.AperiodicOrder.Crystal.Crystals.bead_loadBearing_of_condensed_crosses_threshold, Obsmetrologia.AperiodicOrder.Crystal.Quasicrystals.beadPlacement, Obsmetrologia.AperiodicOrder.Crystal.Quasicrystals.beadCondensedContacts, Obsmetrologia.AperiodicOrder.Crystal.Quasicrystals.bead_loadBearing_of_condensed_crosses_threshold, Obsmetrologia.AperiodicOrder.Crystal.Glass.beadPlacement_eq_crystal_beadPlacement, Obsmetrologia.AperiodicOrder.Crystal.Glass.beadPlacement_eq_quasicrystal_beadPlacement, Obsmetrologia.AperiodicOrder.Crystal.Glass.beadCondensedContacts_eq_quasicrystal_beadCondensedContacts, Obsmetrologia.AperiodicOrder.Crystal.Glass.beadPlacement_loadBearing_of_condensed_crosses_threshold}\leanok\uses{thm:glass-quasicrystal-retains-bead-condensate, thm:percolation-threshold-law, thm:percolation-density-not-sufficient}
The bead layer now carries a finite placement graph, not just a thread count.
Its contact cells are exactly the condensed contacts of the placement, and once
those condensed contacts cross the finite percolation threshold the bead
placement bears load.  Crystals and quasicrystals project that same placement
and threshold test unchanged.  This proves the deterministic count-topology
bridge; it does not prove the frontier critical-fraction law.
\end{theorem}

\begin{theorem}[The bead phase gives the finite crystal spine]\label{thm:glass-bead-phase-crystal-spine}\lean{Obsmetrologia.AperiodicOrder.Crystal.Crystals.bead_phase_spine}\leanok\uses{thm:glass-bead-material-realization, thm:glass-beads-port-law-condensate, thm:glass-bead-multi-adjacency-carrier, thm:glass-bead-placement-percolation}
Under two finite hypotheses---over-occupation of the bead thread and crossing
the retained placement threshold---a crystal keeps the bead-phase spine: legal
thread routing, contradiction-bearing condensate, mirror-stable contradiction
content, the four adjacency readings, the clopen lenticle overlap and its
pulled-back two-cover, and load-bearing placement.  This is what the bead
phase gives crystals before any diffraction layer is added.
\end{theorem}

\begin{theorem}[The bead phase gives the finite quasicrystal spine]\label{thm:glass-bead-phase-quasicrystal-spine}\lean{Obsmetrologia.AperiodicOrder.Crystal.Quasicrystals.bead_phase_spine}\leanok\uses{thm:glass-bead-phase-crystal-spine, thm:glass-quasicrystal-retains-bead-condensate}
Under two finite hypotheses---over-occupation of the bead thread and crossing
the retained placement threshold---a quasicrystal keeps the whole bead-phase
spine: legal thread routing, contradiction-bearing condensate, mirror-stable
contradiction content, the four adjacency readings, the clopen lenticle overlap
and its pulled-back two-cover, load-bearing placement, and the finite universal
diffraction readout.  Compared with the crystal spine, the extra
quasicrystal datum proved is the finite diffraction mass of the crystal occupation
pattern.  It still does not assert analytic diffraction, a critical fraction,
pairwise separation of the four adjacency graphs, or a continuum embedding.
\end{theorem}

\begin{theorem}[The quasicrystal bead phase carries the K4/FOUR analytic spine]\label{thm:glass-quasicrystal-k4-analytic-spine}\lean{Obsmetrologia.AperiodicOrder.Crystal.Quasicrystals.BeadPhaseSpine, Obsmetrologia.AperiodicOrder.Crystal.Quasicrystals.bead_phase_k4_analytic_spine}\leanok\uses{thm:glass-bead-phase-quasicrystal-spine, thm:k4-four-analytic-spine}
Under the same two finite bead-phase hypotheses, a quasicrystal carries the
finite bead spine together with the K4/FOUR analytic spine over its finite
occupation pattern.  This is what the analytic chase proves for crystals and
quasicrystals at the current proof frontier: the material carrier now reaches
the chromatic Rayleigh/gap/diffraction spine without claiming infinite-volume
diffraction, a full spectral theorem, or a continuum crystal/quasicrystal
limit.
\end{theorem}

\begin{proposition}[The bead phase proof boundary]\label{thm:glass-bead-phase-proof-boundary}\lean{Obsmetrologia.AperiodicOrder.Crystal.Crystals.bead_phase_spine, Obsmetrologia.AperiodicOrder.Crystal.Quasicrystals.bead_phase_spine, Obsmetrologia.AperiodicOrder.Crystal.Quasicrystals.bead_phase_k4_analytic_spine}\leanok\uses{thm:glass-bead-phase-crystal-spine, thm:glass-bead-phase-quasicrystal-spine, thm:glass-quasicrystal-k4-analytic-spine, thm:bead-thread-walk-distinct, thm:bead-threading-closed-walk, thm:percolation-critical-fraction, thm:chromatic-obstruction-structural-defect}
The bead phase is finite and named: a bead is an intermediate carrier
between tiles and crystals; it carries port legality, condensate contradiction,
mirror-stable contradiction content, a four-adjacency carrier, a clopen
lenticle overlap, a deterministic placement threshold, the crystal bead spine,
then the quasicrystal finite-diffraction readout, and now the K4/FOUR analytic
spine over that finite occupation pattern.  The remaining claims are frontier,
not consequences of the spine theorem: pairwise separation of
thread/contact/cell/visible adjacency, visible-surface sofic factoring from the
thread walk, graph-threading as a settling closed walk, probabilistic critical
fraction extraction, chromatic obstruction as an irremovable structural defect,
analytic diffraction measure, a full spectral theorem, and continuum embedding.
\end{proposition}

\begin{theorem}[A glass substance descends in its basin landscape]\label{thm:glass-substance-descent-lowers-energy}\lean{Obsmetrologia.AperiodicOrder.Crystal.Glass.descent_lowers_energy, Obsmetrologia.AperiodicOrder.Crystal.Glass.descent_lowers_rationalEnergy, Obsmetrologia.AperiodicOrder.GlassDynamics.BasinGraph.rationalEnergy_le_of_descendChain}\leanok\uses{thm:glass-substance-bundle, thm:basin-no-descending-cycle}
The descent witness carried by a glass substance strictly lowers the retained
basin energy, and the same descent lowers the internal rational free-energy
lift of that basin landscape.
\end{theorem}

\begin{theorem}[A glass substance retains stationary temperature]\label{thm:glass-substance-temperature}\lean{Obsmetrologia.AperiodicOrder.Crystal.Glass.temperature_eq_distribution, Obsmetrologia.AperiodicOrder.Crystal.Glass.produces_creation_landauer_and_discharge}\leanok\uses{thm:glass-substance-bundle, thm:finite-thermodynamic-transition-production}
The temperature value retained by a glass substance is the stationary
distribution at its distinguished state, and the lifted transition carries
positive creation, balanced flux, positive Landauer cost, and positive
hard-flow discharge.
\end{theorem}

\begin{theorem}[Basin thermalization is glass temperature]\label{thm:glass-basin-thermalization}\lean{Obsmetrologia.AperiodicOrder.Crystal.Glass.BasinThermalization.temperature_eq_distribution, Obsmetrologia.AperiodicOrder.Crystal.Glass.BasinThermalization.ledgerAt_eq_eulerChar, Obsmetrologia.AperiodicOrder.Crystal.Glass.BasinThermalization.produces_landauer_and_discharge}\leanok\uses{thm:glass-substance-temperature, thm:anti-landauer-index-balance, thm:markov-stationary-closure-fixed}
A basin-indexed thermalization puts the stochastic kernel directly on the
metastable basin landscape.  Its temperature is the stationary distribution at
the distinguished basin, each basin carries a Landauer ledger equal to the Euler
characteristic of its support, and the distinguished basin carries balanced flux,
positive Landauer cost, and positive hard-flow discharge.
\end{theorem}

\begin{theorem}[Basin matter ledger counts \(H^1\)]\label{thm:glass-basin-h1-landauer-ledger}\lean{Obsmetrologia.AperiodicOrder.Crystal.Glass.BasinThermalization.h1At_eq_landauerBitsAt, Obsmetrologia.AperiodicOrder.Crystal.Glass.BasinThermalization.antiLandauerBitsAt_eq_h0At, Obsmetrologia.AperiodicOrder.Crystal.Glass.BasinThermalization.h0_minus_h1_eq_ledgerAt, Obsmetrologia.AperiodicOrder.Crystal.Glass.BasinThermalization.loadBearingAt_iff_positive_h1At, Obsmetrologia.AperiodicOrder.Crystal.Glass.BasinThermalization.produces_h1_landauer_matter_ledger}\leanok\uses{thm:glass-basin-thermalization, thm:anti-landauer-index-balance}
The basin support is now read as a matter ledger: the \(H^1\) rank of the basin
cover nerve is exactly its Landauer bit count, the \(H^0\) rank is the anti-
Landauer side, the signed ledger is \(H^0-H^1\), and load-bearing is equivalent
to positive \(H^1\).  This is the finite H¹/Landauer spine indexed by glass
basins.
\end{theorem}

\begin{theorem}[Frozen--hopping threshold]\label{thm:glass-frozen-hopping-threshold}\lean{Obsmetrologia.AperiodicOrder.Crystal.Glass.BasinThermalization.frozenAt_iff_distribution_le, Obsmetrologia.AperiodicOrder.Crystal.Glass.BasinThermalization.hoppingAt_iff_le_distribution, Obsmetrologia.AperiodicOrder.Crystal.Glass.BasinThermalization.threshold_total, Obsmetrologia.AperiodicOrder.Crystal.Glass.BasinThermalization.frozen_and_hopping_eq_temperature}\leanok\uses{thm:glass-basin-thermalization, thm:real-stochastic-interval-value-category}
The ordered stochastic carrier of a basin thermalization defines frozen and
hopping predicates relative to a threshold.  Stationarity rewrites both
predicates as comparisons with the basin distribution value.  Any total
instance of that order classifies the basin as frozen or hopping, and
antisymmetry identifies a threshold that is both above and below the temperature
with the temperature itself.
\end{theorem}

\begin{theorem}[Rational basin threshold]\label{thm:glass-rational-basin-threshold}\lean{Obsmetrologia.AperiodicOrder.Crystal.Glass.RatBasinThermalization.temperature_eq_distribution, Obsmetrologia.AperiodicOrder.Crystal.Glass.RatBasinThermalization.frozenAt_iff_distribution_le, Obsmetrologia.AperiodicOrder.Crystal.Glass.RatBasinThermalization.hoppingAt_iff_le_distribution, Obsmetrologia.AperiodicOrder.Crystal.Glass.RatBasinThermalization.threshold_total, Obsmetrologia.AperiodicOrder.Crystal.Glass.RatBasinThermalization.ledgerAt_eq_eulerChar, Obsmetrologia.AperiodicOrder.Crystal.Glass.RatBasinThermalization.h1At_eq_landauerBitsAt, Obsmetrologia.AperiodicOrder.Crystal.Glass.RatBasinThermalization.loadBearingAt_iff_positive_h1At, Obsmetrologia.AperiodicOrder.Crystal.Glass.RatBasinThermalization.produces_h1_landauer_matter_ledger, Obsmetrologia.AperiodicOrder.Crystal.Glass.RatBasinThermalization.produces_landauer_and_discharge}\leanok\uses{thm:glass-frozen-hopping-threshold, thm:glass-basin-h1-landauer-ledger, thm:real-stochastic-interval-value-category}
For the corpus rational stochastic carrier, the threshold dichotomy needs no
extra totality hypothesis: rational order classifies each distinguished basin
as frozen or hopping relative to any rational threshold.  The same specialized
carrier still projects to stationary temperature, the H¹/Landauer matter
ledger, balanced flux, positive Landauer cost, and positive hard-flow discharge.
\end{theorem}

\begin{theorem}[Basin free-energy bridge]\label{thm:glass-basin-free-energy-bridge}\lean{Obsmetrologia.AperiodicOrder.Crystal.Glass.BasinFreeEnergyBridge.orbit_Omega_le_init, Obsmetrologia.AperiodicOrder.Crystal.Glass.BasinFreeEnergyBridge.minimizer_is_equilibrium, Obsmetrologia.AperiodicOrder.Crystal.Glass.BasinFreeEnergyBridge.temperature_eq_distribution, Obsmetrologia.AperiodicOrder.Crystal.Glass.BasinFreeEnergyBridge.ledgerAt_eq_eulerChar, Obsmetrologia.AperiodicOrder.Crystal.Glass.BasinFreeEnergyBridge.produces_temperature_landauer_and_descent}\leanok\uses{thm:glass-basin-thermalization, thm:glass-substance-descent-lowers-energy, thm:positive-temperature-free-energy-flow}
A basin free-energy bridge puts the lower liquid DFT free-energy descent on the
same basin carrier as the stochastic temperature and Landauer ledger.  Along
the basin relaxation orbit, rational free energy stays below its initial value;
minimizers are fixed points; and the same bridge retains stationary temperature,
Euler-characteristic ledger, balanced flux, positive Landauer cost, and positive
hard-flow discharge.
\end{theorem}

\begin{theorem}[Rational glass rung]\label{thm:rational-glass-rung}\lean{Obsmetrologia.AperiodicOrder.Crystal.Glass.RatBasinFreeEnergyBridge.orbit_Omega_le_init, Obsmetrologia.AperiodicOrder.Crystal.Glass.RatBasinFreeEnergyBridge.temperature_eq_distribution, Obsmetrologia.AperiodicOrder.Crystal.Glass.RatBasinFreeEnergyBridge.threshold_total, Obsmetrologia.AperiodicOrder.Crystal.Glass.RatBasinFreeEnergyBridge.ledgerAt_eq_eulerChar, Obsmetrologia.AperiodicOrder.Crystal.Glass.RatBasinFreeEnergyBridge.produces_glass_rung, Obsmetrologia.AperiodicOrder.Crystal.Glass.RatBasinFreeEnergyBridge.produces_h1_glass_rung}\leanok\uses{thm:glass-rational-basin-threshold, thm:glass-basin-free-energy-bridge, thm:glass-basin-h1-landauer-ledger}
The rational glass rung puts rational stochastic temperature, rational
free-energy descent, unconditional frozen-or-hopping thresholding, and the
H¹/Landauer ledger on the same basin carrier.  This is the finite rational
version of the bridge; it still stops short of an analytic real quotient
temperature or critical basin-connectivity transition.
\end{theorem}

\begin{theorem}[Finite basin-connectivity glass threshold]\label{thm:finite-basin-connectivity-glass-threshold}\lean{Obsmetrologia.AperiodicOrder.Crystal.Glass.RatBasinFreeEnergyBridge.BasinConnectivityTransition, Obsmetrologia.AperiodicOrder.Crystal.Glass.RatBasinFreeEnergyBridge.BasinConnectivityTransition.connected_later, Obsmetrologia.AperiodicOrder.Crystal.Glass.RatBasinFreeEnergyBridge.BasinConnectivityTransition.glassy_iff_temperature_below_connectivity_transition, Obsmetrologia.AperiodicOrder.Crystal.Glass.RatBasinFreeEnergyBridge.BasinConnectivityTransition.glassy_iff_distribution_below_connectivity_transition, Obsmetrologia.AperiodicOrder.Crystal.Glass.RatBasinFreeEnergyBridge.produces_connectivity_glass_rung}\leanok\uses{thm:rational-glass-rung, thm:persistent-h0-components-merge}
A finite basin-connectivity transition retains a persistent-\(H^0\) filtration
on the basin graph and a rational threshold.  The retained connection persists
to later filtration stages, and the glassy predicate is exactly temperature
below the retained connectivity threshold, equivalently the stationary basin
distribution is below that threshold.
\end{theorem}

\begin{theorem}[Temperature bridge and glass rung]\label{thm:temperature-bridge-glass-rung}\lean{Obsmetrologia.Thermodynamics.Temperature.ZeroTemperatureFlow.recovers_nat_lyapunov, Obsmetrologia.Thermodynamics.Temperature.zero_temperature_recovers_nat_lyapunov, Obsmetrologia.Thermodynamics.Temperature.balancedTemperature_unique, Obsmetrologia.Thermodynamics.Temperature.rational_temperature_unique, Obsmetrologia.Thermodynamics.Temperature.stationary_temperature_unique, Obsmetrologia.Thermodynamics.Temperature.detailed_balance_temperature_unique, Obsmetrologia.AperiodicOrder.Crystal.Glass.BasinThermalization.temperature_unique, Obsmetrologia.AperiodicOrder.Crystal.Glass.BasinFreeEnergyBridge.temperature_bridge_and_glass_rung, Obsmetrologia.AperiodicOrder.Crystal.Glass.RatBasinFreeEnergyBridge.produces_connectivity_glass_rung}\leanok\uses{thm:markov-stationary-closure-fixed, thm:zero-temperature-recovery, thm:glass-basin-free-energy-bridge, thm:glass-frozen-hopping-threshold, thm:finite-basin-connectivity-glass-threshold}
The finite theorem proved by the current library is the bundled bridge: at a
stationary basin kernel there is a unique balanced temperature value, namely the
common stochastic advance/distribution value; over the rational carrier,
nonzero discharge also gives a unique quotient temperature characterized by
\(\mathrm{discharge}\cdot T=\mathrm{mint}\); the \(T=0\) endpoint recovers the
Nat-valued hard-closure Lyapunov descent; the basin free-energy step is
non-increasing; the bridge carries the H¹/Landauer matter ledger, balanced flux,
positive Landauer and hard-flow discharge evidence; and the finite glass
predicate is exactly the temperature-below-connectivity-threshold comparison.
\end{theorem}

\begin{theorem}[Real ratio temperature and basin-connectivity transition]\label{thm:real-ratio-temperature-connectivity-frontier}\notready\uses{thm:stage-temperature, thm:temperature-bridge-glass-rung, thm:glass-basin-h1-landauer-ledger, thm:real-stochastic-interval-value-category, thm:percolation-threshold-law, thm:reaction-detailed-balance}
The stronger analytic statement remains red: construct a real-valued
temperature \(T=\mathrm{mint}/\mathrm{discharge}\) from independent creation and
free-energy discharge rates; prove a detailed-balance biconditional
\(\mathrm{detailedBalance}\iff T=T_{\mathrm{eq}}\); and extract a critical
basin-connectivity transition from the filtration rather than retaining it as
input data.  The current Lean
theorem proves the finite stationary common-value, the rational nonzero-
discharge quotient form, the zero-temperature Nat endpoint, the H¹/Landauer
basin ledger, and the retained finite connectivity-threshold form, not this
full real quotient, biconditional, analytic limit, or critical-threshold
extraction.
\end{theorem}
